% Part: normal-modal-logic
% Chapter: syntax-and-semantics
% Section: introduction

\documentclass[../../../include/open-logic-section]{subfiles}

\begin{document}

\olfileid{nml}{syn}{int}

\olsection{પરિચય}

મોડલ તર્કશાસ્ત્ર \emph{મોડલ વિધાનો} અને તેમની વચ્ચેના ફલિતતા
સંબંધોનો અભ્યાસ કરે છે. મોડલ વિધાનોના દાખલા આ છે:
\begin{enumerate}
\item $2+2=4$ હોવું અનિવાર્ય છે.
\item આવતી કાલે વરસાદ પડવો શક્ય છે તે અનિવાર્ય છે.
\item જો~$!A$ શક્ય હોવું અનિવાર્ય હોય, તો~$!A$
  શક્ય છે.
\end{enumerate}
શક્યતા અને અનિવાર્યતા જ એકમાત્ર મોડલતાઓ નથી:
બીજા એકસ્થાની સંયોજકોને પણ મોડલતાઓ ગણવામાં આવે છે.
દાખલા તરીકે, ``$!A$ હોવું જોઈએ,'' ``$!A$ થશે,''
``ડાના જાણે છે કે~$!A$,'' અથવા ``ડાના માને છે કે~$!A$.''

મોડલ તર્કશાસ્ત્ર સૌપ્રથમ એરિસ્ટોટલના \emph{De
  Interpretatione}માં દેખાય છે. અનિવાર્યતામાંથી શક્યતા
ફલિત થાય છે, પરંતુ ઊલટું નહીં; શક્યતા અને અનિવાર્યતા
એકબીજાની મદદથી વ્યાખ્યાયિત થઈ શકે છે; જો $!A \land !B$
શક્યપણે સત્ય હોય તો $!A$ અને $!B$ બંને શક્યપણે સત્ય છે,
પરંતુ ઊલટું નહીં; અને જો $!A \to !B$ અનિવાર્ય હોય,
તો $!A$ અનિવાર્ય હોવાથી~$!B$ પણ અનિવાર્ય છે---આ બાબતો
સૌપ્રથમ તેમણે નોંધેલી.

મોડલ તર્કશાસ્ત્ર પ્રત્યેનો પ્રથમ આધુનિક અભિગમ સી.~આઇ.~લુઇસના
કાર્યમાં જોવા મળે છે; તેનો પરિપાક લુઇસ અને લેંગફોર્ડના
\emph{Symbolic Logic} (1932)માં થયો. લુઇસ અને લેંગફોર્ડને
વસ્તુગત શરતી વડે પ્રેરણનું નિરૂપણ સંતોષકારક લાગતું ન હતું:
``$!A$માંથી $!B$ ફલિત થાય છે'' તેના માટે $!A \lif !B$
યોગ્ય વિકલ્પ નથી. તેથી તેમણે પ્રેરણને ``અનિવાર્યપણે,
જો $!A$ તો $!B$'' તરીકે લક્ષણિત કરવાનો પ્રસ્તાવ કર્યો,
જે $!A \strictif !B$ વડે દર્શાવાય છે. જુદા ગુણધર્મોને
વ્યવસ્થિત કરતાં લુઇસે પાંચ જુદી મોડલ પ્રણાલીઓ ઓળખી:
\Log{S1}, \ldots, \Log{S4}, \Log{S5}; છેલ્લી બેનો આજે પણ
ઉપયોગ થાય છે.

લુઇસ અને લેંગફોર્ડનો અભિગમ સંપૂર્ણપણે \emph{વાક્યરચનાત્મક}
હતો: તેમણે યોગ્ય સ્વયંસિદ્ધિઓ અને નિયમો ઓળખ્યાં અને
તેમના વડે શું સાબિત થઈ શકે તેની તપાસ કરી. અર્થવિજ્ઞાનલક્ષી
અભિગમ લાંબા સમય સુધી હાથ ન લાગ્યો. રુડોલ્ફ કાર્નાપે
\emph{Meaning and Necessity} (1947)માં \emph{અવસ્થા-વર્ણન}
વાપરી પ્રથમ પ્રયત્ન કર્યો; એટલે પરમાણુ વાક્યોનો એવો સંગ્રહ,
જે તે અવસ્થા-વર્ણનમાં ``સત્ય'' હોય. સત્યની વ્યાખ્યા કોઈપણ
વાક્ય~$!A$ સુધી વિસ્તારીને કાર્નાપે $!A$ને \emph{અનિવાર્યપણે
સત્ય} ત્યારે કહ્યું જ્યારે તે બધાં અવસ્થા-વર્ણનોમાં સત્ય હોય.
કાર્નાપનો અભિગમ \emph{પુનરાવર્તિત} મોડલતાઓને સંભાળી શકતો
ન હતો: ``શક્યપણે અનિવાર્યપણે \ldots શક્યપણે $!A$''
સ્વરૂપનાં વાક્યો હંમેશા સૌથી અંદરની મોડલતામાં સંકોચાઈ જાય છે.

મોડલ અર્થવિજ્ઞાનમાં મોટો સીમાચિહ્ન સૉલ ક્રિપ્કેનો લેખ
``A Completeness Theorem in Modal Logic'' (JSL 1959) હતો.
ક્રિપ્કેએ લાઇબ્નિત્સના એ વિચારને આધાર બનાવ્યો કે કોઈ વિધાન
``બધાં શક્ય જગતોમાં'' સત્ય હોય તો તે અનિવાર્યપણે સત્ય છે.
પરંતુ આ વિચારને પણ કાર્નાપના અભિગમ જેવી જ મર્યાદા છે:
કોઈ જગત~$w$ (અથવા અવસ્થા-વર્ણન~$s$)માં વિધાનનું સત્ય
$w$ પર જરાય આધાર રાખતું નથી. તેથી ક્રિપ્કેએ માન્યું કે
જગતો વચ્ચે \emph{સુલભતા સંબંધ} $R$ છે, અને
``અનિવાર્યપણે $!A$'' સ્વરૂપનું વિધાન જગત~$w$માં ત્યારે અને
માત્ર ત્યારે જ સત્ય છે જ્યારે $w$માંથી \emph{સુલભ}
બધાં જગતો~$w'$માં $!A$ સત્ય હોય. આ અભિગમનું કોઈ સ્વરૂપ
આપતા અર્થવિજ્ઞાનને ક્રિપ્કે અર્થવિજ્ઞાન કહે છે. તેનાથી
મોડલ તર્કશાસ્ત્રોના ઝડપી વિકાસને વેગ મળ્યો.

ક્રિપ્કે અર્થવિજ્ઞાન મુજબ અર્થઘટન કરતાં મોડલ તર્કશાસ્ત્ર
\emph{સંબંધાત્મક સંરચનાઓ} ``અંદરથી'' કેવી દેખાય છે તે બતાવે છે.
સંબંધાત્મક સંરચના એટલે દ્વિસ્થાની સંબંધ ધરાવતો ગણ
(દાખલા તરીકે, વર્ગના વિદ્યાર્થીઓનો ગણ, જેને તેમના
સામાજિક સુરક્ષા ક્રમાંક મુજબ ક્રમબદ્ધ કર્યો હોય).
પરંતુ સંબંધાત્મક સંરચનાઓ અનેક ક્ષેત્રોમાં મળે છે:
જગતની અવસ્થાઓની સાપેક્ષ શક્યતા ઉપરાંત, કોઈ વ્યક્તિની
જ્ઞાનાત્મક અવસ્થાઓ જ્ઞાનાત્મક શક્યતા વડે જોડાયેલી હોય
અથવા કોઈ ગતિશીલ તંત્રની અવસ્થાઓ અવસ્થા-સંક્રમણ વડે
જોડાયેલી હોય, વગેરે. આ બધાંનું નિદર્શન મોડલ તર્કશાસ્ત્રથી
થઈ શકે છે: પ્રથમથી સામાન્ય સત્યલક્ષી મોડલ તર્કશાસ્ત્ર,
અને બીજાંથી જ્ઞાનમીમાંસક તર્કશાસ્ત્ર, ગતિશીલ તર્કશાસ્ત્ર
વગેરે મળે છે.

અહીં અમે એક વિશિષ્ટ દૃષ્ટિકોણ પર ધ્યાન આપીએ છીએ, જેને
મોડલ તર્કશાસ્ત્રીઓ ``અનુરૂપતા સિદ્ધાંત'' કહે છે. ક્રિપ્કેની
આરંભની મહત્વની શોધોમાંની એક એ છે કે સુલભતા સંબંધ~$R$ના
ઘણા ગુણધર્મો (જેમ કે પરંપરિતતા, સમમિતિ વગેરે) યોગ્ય
``મોડલ યોજનાઓ'' વડે \emph{મોડલ ભાષામાં જ} લક્ષણિત કરી
શકાય છે. દાખલા તરીકે, મોડલ તર્કશાસ્ત્રીઓ કહે છે કે
$R$ની સ્વવાચકતા ``જો અનિવાર્યપણે~$!A$, તો~$!A$''
યોજનાને ``અનુરૂપ'' છે. અમે મુખ્યત્વે \Ax{D}, \Ax{T},
\Ax{B}, \Ax{4} અને~\Ax{5} યોજનાઓના સંયોજનથી મળતી
મોડલ તર્કશાસ્ત્રની કેટલીક શાસ્ત્રીય પ્રણાલીઓ (જેમ કે
\Log{S4} અને \Log{S5})ના અનુરૂપતા સિદ્ધાંતની તપાસ કરીએ છીએ.

\end{document}
